\documentclass[10pt,letterpaper]{article}
\usepackage[letterpaper,top=0.48in,bottom=0.48in,left=0.58in,right=0.58in]{geometry}
\usepackage[T1]{fontenc}
\usepackage{lmodern}
\usepackage{microtype}
\usepackage[hidelinks]{hyperref}
\usepackage{multicol}
\setlength{\parindent}{0pt}
\setlength{\parskip}{2.5pt}
\pagestyle{empty}
\begin{document}
\begin{center}
{\fontsize{12}{13}\selectfont\bfseries CLARA-Care: A Governed Longitudinal Health-AI Prototype Linking Interoperable Evidence, Task-Bounded Context, and Persistent Writeback}\par
\vspace{2pt}
{\bfseries Nguyen Ngoc Thien; Trinh Minh Quang}\par
Hai Ba Trung High School, Hue, Vietnam - Correspondence: \href{mailto:kakangocthien109@gmail.com}{kakangocthien109@gmail.com}
\end{center}
\fontsize{9.1}{9.9}\selectfont

\textbf{Abstract.} CLARA-Care is a research prototype for longitudinal personal-health AI that connects source-preserving data ingestion, task-bounded model context, and governed persistent writeback. The demonstration follows one synthetic record from FHIR-derived evidence through longitudinal state construction, a purpose-bounded AI snapshot, a model-derived proposal, version and governance checks, and an auditable accept-or-reject decision. The system remains a research prototype and is not deployed for clinical care.

\textbf{Purpose and problem.} A longitudinal assistant repeatedly reads from a growing health record and may later propose changes to it. Relevant retrieval alone does not guarantee that the model sees the correct current state, that the disclosed context is appropriate for the present task, or that a proposal remains admissible if the record or governance conditions change before persistence. CLARA-Care makes these boundaries explicit and prevents model output from becoming a direct write to canonical state.

\textbf{System and demonstration.} The demo uses synthetic data and exposes the full decision trace. A version-explicit ingestion component accepts supported FHIR STU3/R4 Bundles, validates subject references, preserves source resources, and separates clinical valid/effective time from the time evidence became known. The longitudinal state layer keeps provenance and version coordinates. A Task-Bounded Health State Snapshot (THSS) then projects only the context admitted for a defined subject, actor, purpose, and task. On the evaluated snapshot-bound path, if the model proposes a durable change, that proposal returns through a Governed State Transition (GST), where state freshness and snapshot binding are checked before persistence; strict research paths also revalidate relevant governance. The interface shows which evidence was selected, which coordinates bound the proposal, and why the write was accepted or rejected. The live demonstration includes a clean control, a stale-state case, and a governance-change case.

\textbf{Evaluation evidence.} In a frozen 12-schedule PostgreSQL governance-TOCTOU matrix, the prototype observed no forbidden commit across consent revocation, role change, policy-epoch advance, state change, and compound drift; invalid post-change proposals were rejected while temporally valid controls remained admissible. A separate bounded model explored 21,361 states and 90,432 transitions to depth 5 with no violations of 11 specified invariants. In an earlier frozen 64-subject synthetic model-context cohort, Strict THSS was all-axis exact for 63/64 Claude subjects and 64/64 Gemini subjects. A newer source-disjoint six-task utility sanity grid completed 48 calls across two locked model families and observed 11/12 correct Strict-THSS cells, 12/12 bound-THSS cells, 11/12 state-only cells, and 9/12 unbound cells. The grid is too small for a powered superiority or noninferiority claim. A separate sealed 384-subject Claude-only comparison between Strict THSS and full authorized history was null, so the model-context evidence is mixed rather than uniformly favorable. These are software and synthetic research results, not clinical validation.

\textbf{Innovation and current deployment status.} This demonstration does not make a separate algorithmic-novelty claim. The underlying GLHS research evaluates a path that links the exact governed health disclosure used during inference to later write admission; the demo shows how that contract operates in an end-to-end, FHIR-facing prototype without claiming mandatory provenance retention across every internal adaptation route. CLARA-Care is a student-led research system under active development and isolated testing. It has not been deployed for patient care, and no claim is made about clinical effectiveness or regulatory compliance.

\textbf{References.} [1] HL7 International, FHIR R4 specification: Consent, Provenance, and RESTful API, 2019. [2] W3C, PROV-O: The PROV Ontology, 2013. [3] X. Li, Y. Wang, H. Lu, et al., ``MemTX: Transactional Belief Commit for Stateful Agent Memory,'' arXiv:2607.23929, 2026. [4] I. Santos-Grueiro, ``Temporary Authority, Permanent Effects: Commit-Time Authorization for LLM Agents,'' arXiv:2607.10487, 2026.
\end{document}
